\documentclass{situationist-preprint}

\title{Essence Without Modal Collapse}
\msid{MET-13-b}
\msfamily{MET-13}
\msdiscipline{Metaphysics}
\mstype{Consequence}
\mskeywords{essence; real definition; modal collapse; contingency; complete concepts; cosmos; laws of nature}

\begin{document}
\maketitle

\begin{abstract}
MET-13-a argues that every object has a real definition and that necessities hold in virtue of essences. It is often objected that such a view leads to modal collapse: if the cosmos, or the actual world, or each individual, has an essence, and essential truths are necessary, then every truth is necessary. We derive from the conclusion of MET-13 that it does not. A truth is contingent exactly when it is not entailed by the totality of essential truths, and collapse requires that every truth be so entailed. Because real definitions are sparse, stating what a thing is and not everything true of it, the premise fails. We examine the four routes to it, through the cosmos, the actual world, complete individual concepts and the lives of persons, and show that each confuses what a thing is with what is true of it.
\end{abstract}

\section{The criterion of contingency}\label{s:criterion}

We use the notation of MET-13-a: $\nec_x\,p$ says that it is true in virtue of the nature of $x$ that $p$, and Principle~P3 of that manuscript says that if $\nec_x\,p$, then $\nec p$.

\begin{definition}[Essential base]\label{d:base}
The \emph{essential base} $\mathcal{E}$ is the set of propositions $p$ such that $\nec_x\,p$ for some object or objects $x$.
\end{definition}

\begin{lemma}\label{l:closure}
Every proposition entailed by $\mathcal{E}$ is necessary.
\end{lemma}

\begin{proof}
Each member of $\mathcal{E}$ is necessary by Principle~P3 of MET-13-a, and what is entailed by necessary propositions is necessary.
\end{proof}

TSM-10 supplies the converse: metaphysical necessity is truth in all worlds consistent with the essences of things, so a proposition is necessary only if every world consistent with $\mathcal{E}$ makes it true, that is, only if $\mathcal{E}$ entails it.

\begin{proposition}[Criterion]\label{p:criterion}
A true proposition is contingent if and only if it is not entailed by $\mathcal{E}$. Hence modal collapse, the thesis that every truth is necessary, holds if and only if every truth is entailed by the essential base.
\end{proposition}

\begin{proof}
From Lemma~\ref{l:closure} and the account of necessity of TSM-10, a proposition is necessary exactly when $\mathcal{E}$ entails it.
\end{proof}

Proposition~\ref{p:criterion} locates the issue. The doctrine of MET-13-a that every object has a real definition says nothing about how much a real definition contains. Collapse requires a further premise, which we call \emph{totality}: every truth is entailed by the real definitions of the things there are. The priority of essence (Theorem~4.2 of MET-13-a) concerns necessary truths only; it does not say that every truth holds in virtue of essences, and the collapse argument needs exactly that stronger claim.

\begin{principle}[Sparseness]\label{p:sparse}
The real definition of an object states its kind and, where it has one, its origin, together with the objects on which it constitutively depends; it does not state the further features the object happens to have.
\end{principle}

Principle~\ref{p:sparse} summarizes the account of Section~5 of MET-13-a: the definition of a substance states its kind and origin, that of a set its members, that of a hole its host. It follows from what a real definition is. A real definition answers the question what a thing is, and Definition~2.2 of MET-13-a characterizes it by the propositions true in virtue of the thing's nature alone; a feature that the thing could lose while remaining what it is cannot belong to that answer without making the definition say more than what the thing is. Sparseness is therefore not an extra restriction but the content of the notion, and it is incompatible with totality as soon as anything has a feature not stated by its kind and origin. We now show that the four arguments offered for totality fail.

\section{Four routes to totality}\label{s:routes}

\paragraph{The cosmos.} Everything that happens is a feature of the cosmos; the cosmos has a real definition; so, it is said, its definition includes everything. The inference is invalid. The definition of the cosmos states what kind of thing it is, the maximal spatiotemporal whole or whatever else it is, and its origin if it has one. MET-09 holds that the cosmos is the only fundamental object and that ordinary objects are abstractions grounded in facts about it. That every fact about a part is grounded in a fact about the whole does not make every fact about the whole part of what the whole is: grounding relates facts, and the facts about the cosmos that ground the existence of oceans are contingent facts about its contents, not essential facts about its nature.

The strongest form of the cosmic route adds origin essentialism. Suppose the cosmos has an initial state and has it essentially, and suppose the laws are deterministic. Then, it seems, every later state is fixed by what the cosmos is. It is not, for the laws are not part of anyone's essence. MET-12 holds that laws are the axioms of the best system summarizing the mosaic, and best-system laws are contingent: the mosaic could have been different, and with it the best system. If the essences of dispositional properties make some connections necessary (MET-11), which such properties are instantiated after the initial state is still contingent, as MET-12-a notes. A world with the same initial state and different laws is consistent with $\mathcal{E}$, and so by Proposition~\ref{p:criterion} the later history is contingent.

\paragraph{The actual world.} On TSM-07, possible worlds are maximal consistent sets of propositions. The actual world $@$ is a set, and the real definition of a set states its members; so it is essential to $@$ that the proposition that there are oceans belongs to it, and by Principle~P3 of MET-13-a this is necessary. That is correct, and harmless. What is necessary is that the proposition belongs to $@$, not that it is true. That $@$ is the actual world, that is, that all its members are true, is not stated in the definition of $@$, which would be the same set had another world been actual. The argument confuses membership with truth.

\paragraph{Complete concepts.} Leibniz held that the concept of an individual substance contains everything that will ever be true of it, so that the predicate of every true proposition is contained in its subject \citep[§§8, 13]{Leibniz1686}. If real definitions were complete concepts, totality would hold. But a complete concept is not a statement of what a thing is. It is a record of everything true of it, and Leibniz himself had to distinguish the hypothetical necessity of Caesar's crossing the Rubicon from the absolute necessity of what holds by the nature of things. The real definitions of MET-13-a are what the second notion requires: Caesar's definition states that he is a human being with a certain origin, and his crossing the Rubicon follows from it no more than it follows from his being human. The underlying motive for totality is the principle of sufficient reason in its Leibnizian form, that every truth has its reason in the nature of its subject. The priority of essence does not require it. Contingent truths have grounds, but their grounds are other contingent truths, as the grounding relation of MET-14 allows, not essences.

\paragraph{Persons.} If a person's real definition stated her life, everything she did would be necessary, and nothing would be left for choice. The real definition of a human person states her kind, a human organism capable of thought (PID-05), and her origin. What her life is for is not part of it. EXI-09-a reads Sartre's \enquote{essence} as a defining project and shows that no such project is given to a human being before she exists (Corollary~3.4 of EXI-09-a); the present point is the modal counterpart of that result. Since a defining project is not part of a real definition, what a person makes of herself is not entailed by $\mathcal{E}$, and by Proposition~\ref{p:criterion} it is contingent.

\section{Objections}\label{s:objections}

\begin{objection}[Dispositional essentialism]
If every fundamental property has a dispositional essence, the laws are necessary \citep{Bird2007}; with an essential initial state, determinism then yields collapse after all.
\end{objection}

\begin{reply}
Three conditions must hold together: necessary laws, an essential initial microstate and determinism. We deny the second, and the first is contested within the present framework. The origin of a substance is the source from which it comes, and the cosmos, which has no external source, has at most an origin in the sense of an initial condition, which is a contingent fact about its contents and not a statement of what kind of thing it is. Even granting the first and the third, the argument shows only that the history is necessary given the initial state, which is the hypothetical necessity of the previous section. Collapse requires that the initial state itself be essential to something, and nothing in MET-13-a implies that it is.
\end{reply}

\begin{objection}[Sparseness is ad hoc]
Principle~\ref{p:sparse} is introduced to block collapse. Without an independent criterion of what belongs to a definition, it begs the question.
\end{objection}

\begin{reply}
The criterion is independent: a feature belongs to an object's definition if it is true in virtue of what the object is, and the test of Fine's cases, on which the singleton's membership belongs to its nature and not to Socrates', shows that the notion discriminates finely \citep{Fine1994}. Sparse accounts of essence have been defended for reasons that have nothing to do with collapse \citep{Wildman2013}. The kind and origin of an object are what fix its conditions of identity, by Principles~P1 and~P2 of MET-13-a; the rest of its features do not, and that is why they are not part of its definition.
\end{reply}

\section{Conclusion}

Modal collapse follows from essentialism only with the premise that every truth is entailed by the essences of things. Real definitions are sparse: they state kinds, origins and constitutive dependence, not everything true of their objects. The arguments for totality, from the cosmos, the actual world, complete concepts and the lives of persons, confuse what a thing is with what is true of it. Essence is prior to modality, and contingency survives.

\begin{thebibliography}{99}\small
\bibitem[Bird(2007)]{Bird2007} Bird, A. (2007). \emph{Nature's Metaphysics: Laws and Properties}. Oxford University Press.
\bibitem[Fine(1994)]{Fine1994} Fine, K. (1994). Essence and modality. \emph{Philosophical Perspectives}, 8, 1--16.
\bibitem[Leibniz(1686)]{Leibniz1686} Leibniz, G.~W. (1686). \emph{Discourse on Metaphysics}.
\bibitem[Wildman(2013)]{Wildman2013} Wildman, N. (2013). Modality, sparsity, and essence. \emph{The Philosophical Quarterly}, 63(253), 760--782.
\end{thebibliography}

\end{document}
